\documentclass[8pt]{article}


\usepackage[T2A]{fontenc}
\usepackage[utf8]{inputenc}
\usepackage[english,russian]{babel}
\usepackage{blindtext}
\usepackage[a4paper, total={8.0in, 10in}]{geometry}

\usepackage[utf8]{inputenc}


\usepackage{amsmath,amssymb,amsfonts,amsthm}
\usepackage{mathrsfs}  
\usepackage{cancel}
\usepackage{leftindex}  

\usepackage{xcolor}
\setcounter{MaxMatrixCols}{20}
% \usepackage{lscape}

\newcommand{\bo}[1] {\mathbf{#1}}
\newcommand{\R} {\mathbb{R}}
\newcommand{\norm}[1] {\left|\left|{#1}\right|\right|}
\newcommand{\T}{^\top}



% \newcommand{\cos}[1] {\text{cos}\left(#1\right)}
% \newcommand{\sin}[1] {\text{cos}\left(#1\right)}



\newtheorem{assumption}{Assumption}
\newtheorem{remark}{Remark}
\newtheorem{theorem}{Theorem}


\definecolor{myhotpink}{RGB}{255, 80, 200}
\definecolor{mywarmpink}{RGB}{255, 60, 160}
\definecolor{mylightpink}{RGB}{255, 80, 200}
\definecolor{mypalepink}{RGB}{255, 200, 230}
\definecolor{mypink}{RGB}{255, 30, 80}
\definecolor{mydarkpink}{RGB}{155, 25, 60}

\definecolor{mypaleblue}{RGB}{240, 240, 255}
\definecolor{mylightblue}{RGB}{120, 150, 255}
\definecolor{myblue}{RGB}{90, 90, 255}
\definecolor{mygblue}{RGB}{70, 110, 240}
\definecolor{mydarkblue}{RGB}{0, 0, 180}
\definecolor{myblackblue}{RGB}{40, 40, 120}

\definecolor{mygreen}{RGB}{0, 200, 0}
\definecolor{mygreen2}{RGB}{245, 255, 230}

\definecolor{mygray}{gray}{0.8}
\definecolor{mydarkgray}{RGB}{80, 80, 160}
\definecolor{mylightgray}{RGB}{160, 160, 160}
\definecolor{mylightgreyred}{RGB}{210, 160, 160}



\definecolor{mydarkred}{RGB}{160, 30, 30}
\definecolor{mylightred}{RGB}{255, 150, 150}
\definecolor{myred}{RGB}{200, 110, 110}
\definecolor{myblackred}{RGB}{120, 40, 40}

\definecolor{myred1}{RGB}{100, 0, 0}
\definecolor{myred2}{RGB}{160, 60, 60}
\definecolor{myred3}{RGB}{220, 90, 120}
\definecolor{myred4}{RGB}{250, 120, 120}
\definecolor{myred5}{RGB}{255, 150, 150}

\definecolor{mygreen}{RGB}{0, 200, 0}
\definecolor{mygreen2}{RGB}{205, 255, 200}
\definecolor{mydarkgreen}{RGB}{0, 100, 0}

\definecolor{mydarkcolor}{RGB}{60, 25, 155}
\definecolor{mylightcolor}{RGB}{130, 180, 250}



\newcommand{\re}[1] {\textcolor{mydarkpink}{#1}}

%%%%%%%%%%%%%%%%%%%%%%%%%%%%%%%%%%%%%%%

\usepackage{tcolorbox}
\tcbuselibrary{skins,breakable}
\usetikzlibrary{shadings,shadows}

\newenvironment{myexampleblock}[1]{%
    \tcolorbox[beamer,%
    noparskip,breakable,
    colback=mylightgreygreen,colframe=myblackgreen,%
    colbacklower=mylime!75!mylightgreygreen,%
    title=#1]}%
    {\endtcolorbox}

\newenvironment{myalertblock}[1]{%
    \tcolorbox[beamer,%
    noparskip,breakable,
    colback=mylightgreyred,colframe=myblackred,%
    colbacklower=myred!75!mylightgreyred,%
    title=#1]}%
    {\endtcolorbox}

\newenvironment{myblock}[1]{%
    \tcolorbox[beamer,%
    noparskip,breakable,
    colback=mylightblue,colframe=myblackblue,%
    colbacklower=myblackblue!75!myblue,%
    title=#1]}%
    {\endtcolorbox}

\title{MyMathTemplate}


\begin{document}

\section*{Домашняя работа 1 }

Дан манипулятор с 7ю степенями свободы. Якобиан, связывающий линейную и угловую скорости схвата ($\bo{v}$, $\omega$) со скоростями в шарнирах манипулятора $\dot{\bo{q}}$ задан как:

\begin{equation}
    \bo{J} = 
    \begin{bmatrix}
    (0.45+n/20)&  0.38&  0.22&  0.05&  0.03&  0.02&  0.01 \\
    0.32&  0.41&  0.28&  0.04&  0.02&  0.01&  0.01 \\
    0.28&  0.35&  0.52&  0.06&  0.03&  0.02&  0.01 \\
    0.12&  0.45&  0.38&  0.82&  0.15&  0.12&  0.08 \\
    0.08&  0.12&  0.25&  0.18&  0.75&  0.22&  0.15 \\
    0.95&  0.08&  0.15&  0.12&  0.08&  0.85&  (0.72-n/20)
    \end{bmatrix}
\end{equation}

где $n$ - номер студента по списку, $\begin{bmatrix}\bo{v} \\ \omega\end{bmatrix} =
\bo{J}\dot{\bo{q}}$.

\begin{itemize}

\item Задание 1. Найдите наименьшее решение (по норме скоростей в шарнирах $\|\dot{\bo{q}} \|$) позволяющее манипулятору достичь желаемой линейной скорости схвата $\bo{v}^*$

\begin{equation}
    \bo{v}^* = 
    \begin{bmatrix}
    1 \\
    -1 \\
    1
    \end{bmatrix}
\end{equation}

\item Задание 2. Определите какую размерность имеет пространство всех решений, доставляющих линейную скорость схвата $\bo{v}^*$.

\item Задание 3. Найдите наименьшее решение (по норме скоростей в шарнирах) позволяющее манипулятору достичь желаемой линейной скорости схвата $\bo{v}^*$, при условии, что скорости в первых трех шарнирах одинаковые.

\item Задание 4. Найдите наименьшее решение (по норме скоростей в шарнирах) позволяющее манипулятору достичь желаемой линейной скорости схвата $\bo{v}^*$, при условии, что скорость в первом шарнире равна единице.

\item Задание 5. Определите, возможно ли найти решение, позволяющее манипулятору достичь желаемой линейной скорости схвата $\bo{v}^*$, и желаемой угловой скорости схвата $\omega^*= 
    \begin{bmatrix}
    0.2 \\
    0.5 \\
    0.1
    \end{bmatrix}$, при условии что скорость в последнем шарнире равна единице, а скорости в первом и втором шарнире равны между собой.

\item Задание 6. Если мы зафиксируем первые 4 шарнира (их скорости равны нулю), будет ли существовать такая линейная скорость схвата $\bo{v}^*$ которую нельзя достичь? Докажите ответ.

\end{itemize}
\end{document}